\section{Finite-\texorpdfstring{$K$}{K} posterior learning}
\label{sec:finite-k}

We formulate the exact all-fail risk under the prescribed inference protocol and obtain a differentiable finite-$K$ surrogate for posterior learning (\cref{sec:problem-setting}).
Fenchel conjugacy introduces a per-problem dual state that makes the surrogate's posterior-dependent term a one-draw expectation (\cref{sec:fenchel}).
We update this state from running failure estimates through a Bregman update, and its tracking gap measures the error from using a lagged dual state in the surrogate objective (\cref{sec:bregman-tracking}).

\subsection{Posterior learning for fixed sampling budgets}
\label{sec:problem-setting}

Many reasoning problems admit multiple solution paths, so a single solver execution can fail even when another succeeds.
Multi-sample inference uses a fixed budget of $K$ executions to increase the chance that the candidate pool contains a successful solution.
We make the candidate-generating distribution trainable by learning a posterior $q_\phi(\theta)$ over solver parameters, regularized toward a prior $p_0(\theta)$.
At inference, we draw $K$ parameter samples from this posterior and run each sampled solver at the prescribed reasoning depth $D$.
The selector associated with the underlying reasoning method then chooses or aggregates their outputs.
We retain that selector and its established training protocol to isolate posterior learning from selector design.
Learning a new selector or co-training the selector and posterior is a complementary research direction.
Our objective targets candidate coverage.
The retained selector determines which candidate the system returns.

For problem $x$, let $F(\theta,x)\in\{0,1\}$ denote the exact failure indicator for solver $\theta$ at depth $D$, with $F=1$ for failure.
For independent draws $\theta_{1:K}\overset{\mathrm{iid}}{\sim}q_\phi$,
\begin{align}
  \Pr(\text{all $K$ solvers fail}\mid x)
  & =
  \mathbb E_{\theta_{1:K}}
  \left[
    \prod_{k=1}^K F(\theta_k,x)
  \right]
  \notag\\
  & =
  \left(
    \mathbb E_{\theta\sim q_\phi}[F(\theta,x)]
  \right)^K.
  \label{eq:binary-all-fail}
\end{align}
The product is the all-fail indicator.
Independence yields the $K$-th power, so minimizing it maximizes Pass@$K$.

This all-fail probability induces the regularized exact target
\begin{equation}
  \min_\phi
  \mathbb E_x
  \left[
    \left(
      \mathbb E_{\theta\sim q_\phi}[F(\theta,x)]
    \right)^K
  \right]
  +
  \tau_{\mathrm{KL}}D_{\mathrm{KL}}(q_\phi\Vert p_0).
  \label{eq:exact-finite-k-objective}
\end{equation}
The data term is the expected all-fail probability, and the KL term regularizes the solver posterior toward $p_0$~\citep{bissiri2016general,lacostejulien2011losscalibrated}.

The discrete indicator $F$ does not provide a pathwise gradient through the solver.
We therefore replace it during training with a differentiable failure score $f(\theta,x)\in[0,1]$ and write $\bar f_\phi(x)\coloneqq\mathbb E_{\theta\sim q_\phi}[f(\theta,x)]$ for its posterior mean.
The resulting differentiable training objective is
\begin{equation}
  \min_\phi
  \mathbb E_x[\bar f_\phi(x)^K]
  +
  \tau_{\mathrm{KL}}D_{\mathrm{KL}}(q_\phi\Vert p_0).
  \label{eq:finite-k-objective}
\end{equation}
For $K=1$, its data term reduces to the expected surrogate failure of one posterior draw.
When $f=F$, the optimization problems in \cref{eq:finite-k-objective,eq:exact-finite-k-objective} coincide; otherwise, \cref{eq:finite-k-objective} serves as a smooth proxy for the exact all-fail risk.

\subsection{One-draw min--max posterior learning}
\label{sec:fenchel}

The product of $K$ independent failure scores is an unbiased estimator of $\bar f_\phi(x)^K$, but computing it requires $K$ solver executions per update.
A one-draw alternative raises one sampled failure score to the $K$-th power.
This quantity estimates $\mathbb E_{\theta\sim q_\phi}[f(\theta,x)^K]$, whereas the finite-$K$ surrogate contains $\bigl(\mathbb E_{\theta\sim q_\phi}[f(\theta,x)]\bigr)^K$.
We therefore seek an exact reformulation of the surrogate in \cref{eq:finite-k-objective} that avoids $K$ solver executions at every update.

The finite-$K$ surrogate has an exact min--max form.
This form makes the posterior-dependent update a one-draw expectation for a fixed dual state.

Define $\psi_K(\lambda)=(K-1)(\lambda/K)^{K/(K-1)}$ on $[0,K]$.
For $K\geq2$ and $r\in[0,1]$,
\begin{equation}
  r^K
  =
  \max_{\lambda\in[0,K]}
  \{\lambda r-\psi_K(\lambda)\},
  \qquad
  \lambda^\star(r)=Kr^{K-1}.
  \label{eq:fenchel-identity}
\end{equation}
Applying \cref{eq:fenchel-identity} to $r=\bar f_\phi(x)$ gives the equivalent saddle objective
\begin{equation}
  \min_\phi
  \max_{\{\lambda_x\in[0,K]\}_x}
  \left\{
    \mathbb E_x
    \left[
      \mathbb E_{\theta\sim q_\phi}[\lambda_x f(\theta,x)]
      -\psi_K(\lambda_x)
    \right]
    +
    \tau_{\mathrm{KL}}D_{\mathrm{KL}}(q_\phi\Vert p_0)
  \right\}.
  \label{eq:saddle-objective}
\end{equation}
For a finite training set, \cref{eq:saddle-objective} requires one scalar dual state $\lambda_x$ for each training problem.

For a fixed current value of $\lambda_x$, its posterior-dependent term is an ordinary expectation over $\theta\sim q_\phi$.
Assume that $q_\phi$ admits a differentiable reparameterization and that differentiation can pass through the expectation.
One posterior draw then gives an unbiased stochastic gradient of the fixed-dual term.
The optimal dual state is $\lambda_x^\star=K\bar f_\phi(x)^{K-1}$.
If $\lambda_x=\lambda_x^\star$, the one-draw posterior gradient equals the finite-$K$ surrogate data gradient.
Otherwise, it remains unbiased for the current fixed-dual objective, whereas it need not equal the gradient of the nonlinear surrogate objective.
The KL regularizer remains part of the posterior update.

\subsection{Bregman update for the dual state}
\label{sec:bregman-tracking}

The optimal dual state depends on $\bar f_\phi(x)$ and changes as the posterior changes.
Re-estimating it from multiple posterior draws at every update would remove the one-draw advantage.
Instead, we update the dual state from the current sampled failure score $f_t=f(\theta_t,x)$.

Let $D_{\psi_K}(\lambda,\mu)=\psi_K(\lambda)-\psi_K(\mu)-\psi_K'(\mu)(\lambda-\mu)$ denote the Bregman divergence generated by $\psi_K$.
We apply the Bregman-proximal ascent step
\begin{equation}
  \lambda_{x,t+1}
  =
  \arg\max_{\lambda\in[0,K]}
  \left\{
    \lambda f_t-\psi_K(\lambda)
    -\eta^{-1}D_{\psi_K}(\lambda,\lambda_{x,t})
  \right\},
  \qquad \eta>0.
  \label{eq:dual-prox}
\end{equation}
Since $\psi_K'(\lambda)=(\lambda/K)^{1/(K-1)}$, this update has a closed form.
With $\beta=1/(1+\eta)$ and $\widehat f_{x,t-1}=\psi_K'(\lambda_{x,t})$,
\begin{align}
  \widehat f_{x,t}
  &=
  \beta\widehat f_{x,t-1}
  +(1-\beta)f_t,
  \notag\\
  \lambda_{x,t+1}
  &=K\widehat f_{x,t}^{K-1}.
  \label{eq:dual-ema}
\end{align}
The closed form reduces dual ascent to a scalar exponential moving average (EMA).
It requires neither an inner optimization loop nor an additional solver execution.
A larger $\beta$ suppresses sampling noise but increases lag when the posterior failure rate changes.

At update $t$, we compute $\lambda_{x,t}=K\widehat f_{x,t-1}^{K-1}$ from the running estimate available before drawing $\theta_t$.
We hold this state fixed during the posterior update.
Afterward, we incorporate $f_t$ into $\widehat f_{x,t}$ and use the result at the next update.
This order fixes the dual state before the current posterior sample, which gives the conditional one-draw gradient property in \cref{sec:fenchel}.

\begin{algorithm}[H]
  \caption{One-draw finite-$K$ posterior learning}
  \label{alg:finite-k}
  \small
  \begin{algorithmic}[1]
    \Require Posterior $q_\phi$, prior $p_0$, pool size $K$, coefficient $\beta$, and failure estimates $\{\widehat f_x\}$
    \For{each training update $t$}
      \State Set and detach $\lambda_{x,t}\gets K\widehat f_{x,t-1}^{K-1}$ for each minibatch problem $x$
      \State Draw $\theta_t\sim q_\phi$ once and compute $f(\theta_t,x)$
      \State Update $\phi$ with $\lambda_{x,t}f(\theta_t,x)$ and the KL regularizer while holding $\lambda_{x,t}$ fixed
      \State Update $\widehat f_{x,t}\gets\beta\widehat f_{x,t-1}+(1-\beta)f(\theta_t,x)$
    \EndFor
  \end{algorithmic}
\end{algorithm}

Direct finite-$K$ training uses $K$ solver executions per update.
The exact min--max reformulation makes the posterior update a one-draw expectation once the current dual state is held fixed.
The Bregman step gives a closed-form EMA update for that state without additional solver executions.
The practical update follows the finite-$K$ data gradient to the extent that $\widehat f_{x,t-1}$ tracks $\bar f_\phi(x)$.

